\documentclass[11pt,a4paper]{article}
% Shared preamble for the English probability and statistics documents.
% Based on the original course preamble. Compile with pdfLaTeX.
\usepackage[top=1.5cm,bottom=2cm,left=2cm,right=2cm]{geometry}
\usepackage[T1]{fontenc}
\usepackage{lmodern}
\usepackage[utf8]{inputenc}
\usepackage[english]{babel}

\usepackage{amsmath}
\usepackage{amssymb}
\usepackage{amsthm}
\usepackage{mathrsfs}
\usepackage{mathtools}
\usepackage{etoolbox}
\usepackage{tikz}
\usetikzlibrary{shapes.multipart,arrows.meta}
\usepackage{xcolor}
\usepackage{enumerate}
\usepackage[shortlabels]{enumitem}
\usepackage{hyperref}
\usepackage{version}
\usepackage{bbold}
\usepackage{fontawesome5}
\usepackage{fancyhdr}
\lfoot{}
\usepackage{natbib}
\usepackage{graphicx}
\usepackage{pdfpages}
\usepackage{authblk}
    
\usepackage{blindtext}
\usepackage{multicol}
\usepackage{collcell}
\usepackage{wrapfig}
\usepackage{blkarray}
\usepackage{cancel}
\usepackage{appendix}
\usepackage{float}
\makeatletter
\input{tex/assets/slashbox.sty}
\makeatother
\hypersetup{unicode,colorlinks,citecolor=red,filecolor=black,
  linkcolor=red,urlcolor=cyan}
\setlength{\columnsep}{1cm}
\pagenumbering{gobble}

\def\numtext#1{%
  \ifcase#1\or one\or two\or three\or four\or five\or six\or seven\or eight\or nine\or ten\or
  eleven\or twelve\or thirteen\or fourteen\or fifteen\or sixteen\or seventeen\or eighteen\or nineteen\or twenty\or Lots\fi}

\newcommand{\Byte}[1]{%
  \def\nodecontents{}%
  \foreach[count=\l] \k in {#1} {\xappto{\nodecontents}{\noexpand\nodepart{\numtext{\l}}\k}}%
  \node[rectangle split,rectangle split parts=8,rectangle split horizontal,draw] at (2,2){\nodecontents};}

\newcommand{\D}{{\mathrm{d}}}
\newcommand{\cE}{{\mathcal{E}}}
\newcommand{\cG}{{\mathcal{G}}}
\newcommand{\cX}{{\mathcal{X}}}
\newcommand{\cY}{{\mathcal{Y}}}
\newcommand{\cB}{{\mathcal{B}}}
\newcommand{\cF}{{\mathcal{F}}}
\newcommand{\cO}{{\mathcal{O}}}
\newcommand{\cS}{{\mathcal{S}}}
\newcommand{\cT}{{\mathcal{T}}}
\newcommand{\given}{\, \big| \,}
\newcommand{\cL}{{\mathscr{L}}}
\newcommand{\cD}{{\mathcal{D}}}
\newcommand{\cA}{{\mathscr{A}}}
\newcommand{\cP}{{\mathscr{P}}}
\newcommand{\cR}{{\mathscr{R}}}
\newcommand{\A}{{\mathcal A}}
\newcommand{\one}{\mathbb{1}}
\newcommand{\mathbbm}[1]{\mathbb{#1}}

\providecommand{\C}{}
\renewcommand{\C}{\mathbb C}
\newcommand{\R}{\mathbb R}
\newcommand{\Q}{\mathbb Q}
\newcommand{\N}{\mathbb N}
\newcommand{\T}{\mathbb T}
\newcommand{\F}{\mathbb F}
\renewcommand{\P}{\mathbb P}
\newcommand{\Z}{\mathbb Z}
\newcommand{\E}{\mathbb {E}}
\newcommand{\Pbb}[1]{\P\left( #1\right)}
\newcommand{\Ebb}[1]{\E\left[ #1\right]}

\newcommand{\lb}{\left (}
\newcommand{\rb}{\right )}
\newcommand{\lbb}{\left [}
\newcommand{\rbb}{\right ]}
\newcommand{\labs}{\left |}
\newcommand{\rabs}{\right |}
\newcommand{\lbrb}[1]{\lb #1 \rb}
\newcommand{\lbbrbb}[1]{\lbb#1\rbb}
\newcommand{\labsrabs}[1]{\labs#1\rabs}
\newcommand{\lbbrb}[1]{\lbb#1\rb}
\newcommand{\lbrbb}[1]{\lb#1\rbb}
\newcommand{\langlerangle}[1]{\langle#1\rangle}
\newcommand{\lbcurly}{\{}
\newcommand{\rbcurly}{\}}
\newcommand{\curly}[1]{\lbcurly#1\rbcurly}
\newcommand{\abs}{\labsrabs}
\newcommand{\ind}[1]{1_{#1}}
\newcommand{\indep}{\perp \!\!\! \perp}
\renewcommand{\d}{\, \mathrm{d}}

\newtheorem{theorem}{Theorem}
\newtheorem{problemCnt}{Lemma}
\theoremstyle{definition}
\newtheorem{problem}[problemCnt]{Problem}
\newtheorem{problembook}[problemCnt]{\faBook \,  Problem}
\newtheorem{problempencil}[problemCnt]{\faPencil* \,  Problem}
\newtheorem{problemstar}[problemCnt]{\faStar\,  Problem}
\newtheorem{problemstarpencil}[problemCnt]{\faStar \, \faPencil* \,   Problem}
\newtheorem{probleminteresting}[problemCnt]{\faSmile[regular] \,   Problem}
\newtheorem{solution}{Solution}

\newcommand{\fwcell}[1]{\makebox[\arraystretch\normalbaselineskip]{$#1$}}
\newenvironment{enumeratei}{\begin{enumerate}[(i)]}{\end{enumerate}}
\newenvironment{enumeratea}{\begin{enumerate}[(a)]}{\end{enumerate}}
\newenvironment{enumeraten}{\begin{enumerate}[\upshape 1.]}{\end{enumerate}}

% Compatibility helpers needed by a few files in the combined archive.
\AtBeginDocument{\noindent{\footnotesize English translation prepared with ChatGPT (OpenAI, GPT-6), 26 October 2026, for translation of the Bulgarian text into English and LaTeX formatting.
The Bulgarian originals remain the reference; translation may contain errors.}\par\smallskip}
\newcommand{\MissingGraphic}[1]{%
  \begin{center}\fbox{\parbox{0.82\linewidth}{\centering Missing source image: \texttt{#1}}}\end{center}}



\pagestyle{fancy}\fancyhf{}\lhead{PTMS, 2023/2024}\rhead{Homework 2}
\setlength{\headheight}{15pt}
\begin{document}
\emph{The maximum total for all homework assignments during the semester is 50 points,
corresponding to a bonus of 0.5 towards the exercise-class grade.}
\begin{problem}
(3 points) Let $X,Y$ be random numbers from $\{1,\ldots,n\}$\footnote{That is,
independent and uniformly distributed on this set.}, and let $p_n$ be the
probability that $X+Y$ is a perfect square. Find
\[\lim_{n\to\infty}\sqrt n p_n.\]
\end{problem}
\begin{problem}
($6\times2=12$ points)
\begin{enumerate}
\item One card is drawn from each of 10 standard decks of 52 cards. Find the
probability that the resulting 10-card hand:
\begin{enumeratea}
\item contains no repeated cards;
\item contains at least three aces;
\item contains four clubs, three diamonds, two hearts and one spade;
\item contains exactly 4 more black cards than red cards, given that there are more black than red cards.
\end{enumeratea}
\item All items in one batch are functional; in another batch, $1/4$ are defective.
An item from a randomly selected batch turns out to be functional. Calculate the
probability that a second randomly chosen item from the same batch is defective,
if the first item was returned to its batch after inspection.
\item What is the probability that the roots of $x^2+ax+b=0$, $a,b\in[0,1]$, are real?
\item Two faces each of a fair die are coloured white, green and red. Roll it twice.
Let $X$ count white outcomes and $Y$ count red outcomes. Find their joint
distribution, determine whether they are independent, and find their covariance,
$\P(X=1\mid Y=1)$ and $\P(X>Y)$.
\item Let $X$ take values in $\mathbb N_0$, and let
\[g_X(s):=\E s^X=\sum_{k=0}^\infty s^k\P(X=k)\]
be its generating function.
\begin{enumeratea}
\item Let $Y:=3X$ and $Z=X_1+X_2$, where $X_1,X_2\sim X$ are independent.
Express $g_Y$ and $g_Z$ in terms of $g_X$.
\item Let $X\sim Ge(p)$, meaning $\P(X=k)=p(1-p)^k$ for $k\in\mathbb N_0$.
Calculate $g_X$ and use it to find $\E(X)$ and $D(X)$.
\end{enumeratea}
\item Let $X,Y$ be independent random variables with
\[\P(X=-1)=\P(X=1)=\P(Y=1)=2\P(Y=3)=2\P(Y=5)=\frac12.\]
Let $Z_1:=2X+Y+1$, $Z_2:=XY$, and $Z_3=X^Y$. Find their expectations and variances.
\end{enumerate}
\end{problem}
\begin{problem}[5 points]
Let $n\ge3$ and $p$ be natural numbers. Consider the following game with players $A_1,\ldots,A_n$:
\begin{itemize}
\item Each player writes predictions for $(2p+1)$ tosses of a standard coin
(probability $1/2$ for each of heads and tails), using H for heads and T for tails.
\item The coin is tossed $(2p+1)$ times, and the players with the most correct
predictions split a prize of $n!$ leva.
\end{itemize}
Example: $n=4$, $p=1$, predictions $HTH,HHH,TTT,HHT$. With outcome $TTH$,
players 1 and 3 each predict 2 correctly and split $4!$, receiving 12 leva each.
For $n=4$, $p=1$, the same predictions and outcome $TTT$, player 3 predicts all
3 correctly and wins $4!=24$ leva.
For $i\in[n]$, define
\begin{itemize}
\item $X_i=$ ``the number of outcomes correctly predicted by $A_i$'';
\item $G_i=$ ``player $A_i$'s winnings'';
\item $q_k=\P(X_i=k)$, $r_k=\P(X_i\le k)$.
\end{itemize}
Our aim is to calculate player $A_1$'s expected winnings, $\E(G_1)$, under two strategies.

\textbf{Part 1. Preliminary considerations.}
\begin{enumerate}
\item * As usual, we work on a probability space $(\Omega,\mathcal F,\P)$.
What set would you propose for $\Omega$?
\item What is the distribution of $X_i$?
\item Calculate $r_p$.
\item Recall that $X(\Omega)$ denotes all images of elements of $\Omega$ under $X$.
For example, $X_1(\Omega)=\{0,1,\ldots,2p+1\}$, since $A_1$ can predict between
0 and $2p+1$ outcomes correctly. If $\P(A)>0$, define conditional expectation by
\[\E(X\mid A)=\sum_{k\in X(\Omega)}k\P(X=k\mid A).\]
Prove that if $B_1,\ldots,B_s$ form a partition into events of nonzero probability, then
\[\E X=\sum_{k=1}^s\E(X\mid B_k)\P(B_k).\]
\emph{(Expand $\E X$ by definition, apply the law of total probability, and interchange the sums.)}
\end{enumerate}
\textbf{Part 2. All players play randomly.} In this part, assume that all players
choose heads or tails randomly with probability $1/2$, independently of each other.
\begin{enumerate}
\item What is $G_1(\Omega)$, the set of possible winnings for $A_1$?
\item For every $k\in X_1(\Omega)$ and $j\in[n]$, prove that
\[\P\left(G_1=\frac{n!}{j}\,\middle|\,X_1=k\right)
=\binom{n-1}{j-1}(r_k-r_{k-1})^{j-1}r_{k-1}^{n-j}.\]
\item Prove that, for $k>0$,
\[\E(G_1\mid X_1=k)=\frac{(n-1)!(r_k^n-r_{k-1}^n)}{q_k}.\]
Note that the formula also holds for $k=0$ if we define $r_{-1}=0$.
\emph{(Expand using the preceding part and use the binomial theorem.)}
\item Conclude that $\E(G_1)=(n-1)!$.
\emph{(Consider the partition $\{X_1=k\}$, $k\in\{0,1,\ldots,2p+1\}$.)}
\item Is the conclusion of the preceding part expected? We could also have
obtained it differently: can we calculate $G_1+\cdots+G_n$? Use linearity of
expectation and symmetry to obtain $\E(G_1)$ again.
\end{enumerate}
\textbf{Part 3. $A_1$ and $A_2$ cooperate; all other players play randomly.}
Assume that everyone except $A_2$ plays randomly. $A_2$ knows $A_1$ and predicts
the opposite of $A_1$; for example, $HHT$ for $A_1$ means $TTH$ for $A_2$.
If either wins, they share the winnings equally. Write
\begin{itemize}
\item $G',G'_1,G'_2$ for their combined and individual winnings, respectively,
so $G'=G'_1+G'_2$;
\item $Y=$ ``the larger number of correct predictions by $A_1$ and $A_2$''
(for predictions $HHT,TTH$ and outcome $HHH$, $Y=2$);
\item as before, for $k\in\{1,3,4,\ldots,n\}$, $q_k=\P(X_i=k)$ and $r_k=\P(X_i\le k)$.
\end{itemize}
\begin{enumerate}
\item Prove that $Y(\Omega)=\{p+1,\ldots,2p+1\}$.
\item Prove that for $k\in\{p+1,\ldots,2p+1\}$, $\P(Y=k)=2q_k$.
\item Prove that for $k\in\{p+1,\ldots,2p+1\}$ and $j\in\{0,\ldots,n-2\}$,
\[\P\left(G'=\frac{n!}{j+1}\,\middle|\,Y=k\right)
=\binom{n-2}{j}(r_k-r_{k-1})^j r_{k-1}^{n-2-j},\]
and consequently, using the binomial theorem,
\[\E(G'\mid Y=k)=n(n-2)!\frac{r_k^{n-1}-r_{k-1}^{n-1}}{q_k}.\]
\item Conclude that
\[\E(G')=2n(n-2)!\left(1-\frac1{2^{n-1}}\right),\]
and check that this strategy is better for $A_1,A_2$ than the strategy in Part 2.
\item Calculate $\E(G_i)$ for $i\in\{3,\ldots,n\}$.
\end{enumerate}
\end{problem}
\end{document}
